% Specimen: Greek report in two columns. Uses every LabForge command.
\documentclass[preset=report, language=greek, bibliography]{.labforge}

\newcommand{\filler}{Ο θάλαμος του τόξου αποτελεί την καρδιά του αντιδραστήρα
και το πρώτο στάδιο της διεργασίας. Σκοπός του είναι να φιλοξενήσει το ηλεκτρικό
τόξο και να παράσχει το περιβάλλον εντός του οποίου πραγματοποιείται η
αποσύνθεση του αζώτου και ο σχηματισμός του μονοξειδίου. Η σχεδίασή του καθόρισε
σε μεγάλο βαθμό τη συνολική γεωμετρία της διάταξης.}

\begin{document}

\myreport[
    header      = ΕΡΓΑΣΤΗΡΙΑΚΗ ΑΝΑΦΟΡΑ,
    title       = Δέσμευση ατμοσφαιρικού αζώτου με ηλεκτρικό τόξο,
    author      = Ομάδα Β,
    conducted   = \mytoday,
    institution = Τμήμα Χημικών Μηχανικών,
]

\begin{multicols}{2}

\mysection[label=intro]{Εισαγωγή}
\filler{} Η μέθοδος βασίζεται στην εργασία των Birkeland και Eyde \cite{birkeland}.

\filler{} Οι θερμοδυναμικές σχέσεις δίνονται στο \cite{atkins}.

\mysubsection[label=chamber]{Θάλαμος τόξου}
\filler{} Η διάταξη φαίνεται στο \myref{photo}, ενώ τα μέρη της στα
\myref{photo, fig:example-image-b}.

\myfigure[label=photo, caption=Ο θάλαμος πριν από τη συναρμολόγηση., width=0.8]{example-image-a.png}
\myfigure[caption=Στερέωση της κάτω φλάντζας σε ξύλινη βάση., width=0.6, place=float]{example-image-b.png}

\filler{} Ο σχηματισμός του μονοξειδίου ακολουθεί την \myref{formation}, με σταθερά
ισορροπίας κατά την \myref{equilibrium}.

\myreaction[label=formation]{N2 + O2 <=> 2 NO}
\myformula[label=equilibrium]{K_p = \frac{p_{\mathrm{NO}}^2}{p_{\mathrm{N_2}}\, p_{\mathrm{O_2}}}}
\myequation{\Delta G^\circ = -RT \ln K_p}

\filler{} Οι μετρήσεις συνοψίζονται στον \myref{tab:yields}.

\mytable[caption={Μάζα, χρόνος αντίδρασης και απόδοση ανά δείγμα.}]{yields.tex}

Ο \myref[nom]{tab:yields} συνεχίζεται στην επόμενη στήλη.
Μικτές αναφορές: \myref{tab:yields, formation}. Ονομαστική: \myref[nom]{tab:yields, formation}.

\mysection[label=results]{Αποτελέσματα}
\filler{} Η καμπύλη βαθμονόμησης φαίνεται στο \myref{plot:absorbance}.

\myplot[caption=Απορρόφηση ως συνάρτηση της συγκέντρωσης., width=0.95]{absorbance.tex}

\filler{} Οι σταθερές του \myref{tab:constants} χρησιμοποιήθηκαν στην \myref*{equilibrium}.

\mytable[caption=Φυσικές σταθερές., width=0.7]{constants.tex}

\myfigure[place=span, caption=Εικόνα σε όλο το πλάτος της σελίδας., width=0.5]{example-image-a.png}

\filler{} Η \myref{results} συνοψίζει τις \myref{intro, chamber}.

\filler{}

\mybibliography
\end{multicols}

\end{document}
